\section{Discussion}
The results establish agreement between implemented contracts and their development oracles. D2 shows why end-to-end graph counts alone are insufficient: repeated baseline structure can dominate totals, while hard-negative source locations expose method-name collisions. Its six corrected v 0.1 errors support targeted diagnosis and regression protection, not an unbiased estimate of improvement on unseen programs.

The three-way decision separates forbidden relationships from changes requiring a design decision. P3 reports historical findings while gating only introduced findings, permitting gradual adoption without silently accepting new violations. REVIEW neither approves evolution nor rewrites the expected model. The evaluated implementation has no durable rejection disposition: rejecting an evolution requires changing or withdrawing the proposed code and rerunning the check; unchanged inputs still yield REVIEW. Accepting it requires a separately reviewed model change. Owner protection is therefore an organizational prerequisite, not an experimentally validated property of the gate.

Component scoping parsed 57/189 file instances and warm caching reduced the recorded median from 518.51 to 242.62ms. Those observations combine parsing, Git processes, cache I/O, comparison, and reporting; they do not isolate algorithmic cost or establish scale. Next experiments require independently annotated real repositories and histories, capability-aligned external baselines, and a practitioner pilot measuring false-block burden, review effort, and model-change governance. IaC, runtime, AI assistance, and other language analyzers remain unevaluated extensions.
